\section{Alternations}


\begin{frame}{Alternations}
\begin{itemize}
\item  Many verbs have multiple possible mappings of grammatical function to role
  \begin{exe}
    \ex
    \begin{xlist}
      \ex\eng{Kim broke the window with the hammer}
      \ex\eng{The hammer broke the window}
      \ex\eng{The window broke}
    \end{xlist}
    \ex
    \begin{xlist}
      \ex\eng{I cut the cake with the knife}
      \ex\eng{This cake cuts easily}
    \end{xlist}
  \end{exe}  
\item  The relations between them are called \txx{alternations}
\item  \textit{English Verb Classes and Alternations} \citep{Levin:1993}
\end{itemize}


\end{frame}

\begin{frame}[allowframebreaks]{There are many alternations}
%\MyLogo{These are also alternations for Levin}
\begin{itemize}
\item  A common  way to change the number of arguments is called 
\txx{voice}: passive, middle
\begin{exe}
  \ex \txx{Transitive Passive}
  \begin{xlist}
    \ex\eng{Kim ate Sandy}
    \ex\eng{Sandy was eaten (by Kim)}
  \end{xlist}
  \ex \txx{Ditransitive Passive} 
  \begin{xlist}
    \ex\eng{Abraham gave Brown chocolate}
    \ex\eng{Abraham gave chocolate to Brown}
    \ex\eng{Chocoloate was given to Brown (by Abraham)}
    \ex\eng{Brown was given chocolate (by Abraham)}
  \end{xlist}
\framebreak
  \ex \txx{Transitive Middle} (or just causative/inchoative)
  \begin{xlist}
    \ex\eng{They open the gate very quietly}
    \ex\eng{The gate opens very quietly}
  \end{xlist}
  \ex \txx{Intransitive Middle}
  \begin{xlist}
    \ex\eng{The knife cuts the cake well}
    \ex\eng{The knife cuts well}
  \end{xlist}
\end{exe}
\item And some more alternations:
\begin{exe}
\ex \txx{Conative} alternation:
\begin{xlist}
    \ex\eng{Kim \ul{hit} the door} $\leftrightarrow$ \eng{Kim \ul{hit} at the door}
  \end{xlist}
\ex \txx{Body-part possessor ascension} alternation:
  \begin{xlist}
    \ex\eng{Kim \ul{cut} Sandy's arm} $\leftrightarrow$ \eng{Kim \ul{cut} Sandy on the arm}
  \end{xlist}

\end{exe}
%\item Can you find 
\end{itemize}

\end{frame}

\begin{frame}{Why so many possibilities?}

\begin{itemize}
\item So we can emphasize different participants
\item We may not know all the participants
\item We may not care about all the participants
\item There are also lexical alternations
\end{itemize}
\begin{exe}
\ex \eng{Kim \ul{killed} Sandy} vs \eng{Sandy \ul{dies}}
\ex c.f. \eng{Kim \ul{melted} the ice} vs \eng{the ice \ul{melted}}
\ex \glll 金が 氷を \ul{溶かした}  {~~~vs~~~}  氷が \ul{溶けた} \\ 
Kim-ga koori-wo \ul{tokashita}  {} koori-ga \ul{toketa} \\
  Kim-\textsc{sbj} ice-\textsc{obj} \textit{melt:trans} 
{} ice-\textsc{sbj} \textit{melt:intrans} \\
\end{exe}


\end{frame}


